\origpage{335}{N. Robertson, P.D. Seymour / Journal of Combinatorial Theory, Series B 90 (2004) 325--385}
\sectionhead{5. Surface homeomorphisms}

Let $\Phi$ be a frame in $\Sigma$. A homeomorphism $\alpha \colon \Sigma \to \Sigma$ is $\Phi$-\textit{preserving} if

\hspace*{12pt}$\bullet$ $\alpha(v) = v$ for all $v \in V(\Phi)$, and $\alpha(S) = S$ for all $S \in E(\Phi)$ (and hence $\alpha$ preserves the direction of each side, and maps each cuff onto itself in an orientation-preserving way)

\hspace*{12pt}$\bullet$ if $\Sigma$ is orientable then $\alpha$ preserves the orientation of $\Sigma$.

Let $(\Gamma, \phi)$ be a $\chi$-coloured painting, and let $\alpha \colon \Sigma_\chi \to \Sigma_\chi$ be a $\Phi_\chi$-preserving homeomorphism. We define the \textit{image} of $(\Gamma, \phi)$ under $\alpha$ in the natural way; that is, it is the $\chi$-coloured painting $(\Gamma', \phi')$ where $U(\Gamma') = \alpha(U(\Gamma))$, $V(\Gamma') = \alpha(V(\Gamma))$, and for each $e \in E(\Gamma), \gamma_{\Gamma'}(\alpha(e)) = \alpha(\gamma_\Gamma(e))$ and $\phi'(\alpha(e)) = \phi(e)$. The proof of the following lemma is clear.

\textbf{5.1.} \textit{Let $(\Gamma_1, \phi_1), (\Gamma_2, \phi_2)$ be $\chi$-coloured paintings, and let $\alpha_1, \alpha_2$ be $\Phi_\chi$-preserving homeomorphisms of $\Sigma$. For $i = 1, 2$ let $(\Gamma'_i, \phi'_i)$ be the image of $(\Gamma_i, \phi_i)$ under $\alpha_i$. Then there is a (linear) inflation of $(\Gamma'_1, \phi'_1)$ in $(\Gamma'_2, \phi'_2)$ if and only if there is one of $(\Gamma_1, \phi_1)$ in $(\Gamma_2, \phi_2)$.}

Thus, for example, if we suppose that $(\Gamma_i, \phi_i)$ $(i = 1, 2, \dots)$ is a bad sequence for $\chi$, then we may replace each $(\Gamma_i, \phi_i)$ by its image under some $\Phi_\chi$-preserving homeomorphism, and obtain another bad sequence; and if the terms of the first sequence are similarly oriented, then so are the terms of the second. This method will be used in combination with the following lemma. Let $\Phi$ be a frame in $\Sigma$. A drawing $K$ in $\Sigma$ is a \textit{feature} (in $\Phi$) if $K$ has no border edges and no isolated vertices, and $U(K) \cap S = \emptyset$ for every short side $S$ of $\Phi$. ($U(K)$ may intersect the long sides of $\Phi$.) Two features $K_1, K_2$ are \textit{equivalent} if there is a $\Phi$-preserving homeomorphism $\alpha$ of $\Sigma$ such that $\alpha$ maps $K_1$ to $K_2$. Hence if $K_1, K_2$ are equivalent they have the same size. (The \textit{size} of a drawing $G$ is $|V(G)| + |E(G)|$.) This defines an equivalence relation, and it follows easily (by choosing a line to represent an edge of $K$, cutting the surface along it, and using induction on $|E(K)|$---see Section 3 of [2] for similar proofs) that

\textbf{5.2.} \textit{For every surface $\Sigma$, every frame $\Phi$ in $\Sigma$ and every integer $n \ge 0$, there are only finitely many equivalence classes of features in $\Phi$ with size $\le n$.}

\fix{We give the proof. Let $K$ be a feature in $\Phi$ with size $\le n$. First we replace $K$ and $bd(\Sigma)$ by a finite family of lines in $\Sigma$, any two of which meet only in common ends. The closure of each edge of $K$ is a line or an $O$-arc; it gives one line of the family in the first case, and two in the second, by splitting it at a point of the edge. The vertices of $\Phi$ and of $K$, together with a finite set of points of $bd(\Sigma)$ that depends only on $\Sigma$, cut $bd(\Sigma)$ into further lines of the family. Each of these lies in the closure of one side of $\Phi$, since its interior contains no vertex of $\Phi$. Thus the number of lines is at most a number $N$ that depends only on $\Sigma$, $\Phi$ and $n$. We number the lines $1, 2, \dots$, and we cut $\Sigma$ along all of them. This gives a surface $Y$, perhaps not connected, with a map onto $\Sigma$ that is one-to-one off $bd(Y)$. Each line of the family not included in $bd(\Sigma)$ has two lifts in $bd(Y)$, and each line in $bd(\Sigma)$ has one, and these lifts cover $bd(Y)$ and meet only in common ends. Let $W$ be a component of $Y$, with $a$ handles, $b$ crosscaps and $c$ cuffs. By the classification of surfaces, $W$ together with the lifts in $bd(W)$ is determined up to homeomorphism by $a$, $b$, $c$ and the cyclic order of the lifts around each cuff of $W$; we call these data the \textit{form} of $W$, and we call the set of the forms of the components of $Y$ the \textit{code} of the family. Here $2a + b$ is at most the corresponding number for $\hat{\Sigma}$, because each choice of values in $\mathbb{Z}_2$ on the $2a + b$ standard loops of $W$ gives a class in $\mathbb{Z}_2$ on $\Sigma$ (since the map from $W$ to $\Sigma$ is one-to-one off $bd(W)$), and different choices give different classes. Also $c$ is at most the number of lifts. Hence for each number of lines there are only finitely many codes. If two families with the same number of lines have the same code, then there is a homeomorphism of $\Sigma$ that maps each line of the first family onto the line of the second with the same number, end to end; for the components with the same form are homeomorphic by maps that keep the lifts, and these maps agree over each line, so together they give a homeomorphism of $\Sigma$.

To each line we attach a \textit{tag}, which says which side of $\Phi$ includes it (if it lies in $bd(\Sigma)$), which of its ends are vertices of $\Phi$, and which ones, and which of its ends are vertices of $K$. Now let $K, K'$ be features with size $\le n$ whose families have the same number of lines, the same code and the same tags, and let $\alpha$ be the homeomorphism just described. Then $\alpha(v) = v$ for every $v \in V(\Phi)$, since each such $v$ is an end of a line, and the tags say that the corresponding end of the corresponding line is $v$. Hence $\alpha(S) = S$ for every side $S$ of $\Phi$, since the closure of $S$ is the union of the lines whose tags assign them to $S$. In the same way $\alpha(U(K)) = U(K')$ and $\alpha(V(K)) = V(K')$. If $\Sigma$ is not orientable, or $\alpha$ preserves the orientation of $\Sigma$, it follows that $K, K'$ are equivalent. Otherwise $\Sigma$ is orientable and $\alpha$ reverses its orientation, since a homeomorphism of a connected orientable surface either preserves or reverses its orientation. If $K', K''$ are both obtained from $K$ in this way, by $\alpha'$ and $\alpha''$ say, then $\alpha'' \circ \alpha'^{-1}$ preserves the orientation, and so $K', K''$ are equivalent. Thus the features with given number of lines, code and tags lie in at most two equivalence classes. Since there are only finitely many choices of the number of lines (at most $N$), of the code and of the tags, this proves 5.2.}\fixnote{fix:5.2-proof-added}{The paper only sketches the proof of 5.2, by an induction on $|E(K)|$. The argument above cuts along all the edges at once; it is the one used in the formalization.}

Now let $\Phi$ be a frame in $\Sigma$, and let $K$ be a feature in $\Phi$. If we cut $\Sigma$ along $U(K)$ we obtain one or several components $\Sigma_1, \dots, \Sigma_k$ say, which we call \textit{fragments} of $\Sigma$ after cutting along $U(K)$. There is a natural surjection of $\psi \colon \Sigma_1 \cup \cdots \cup \Sigma_k \to \Sigma$, which we shall call the \textit{associated surjection}. Let $\Sigma_1$ be a fragment. We denote by $\psi^{-1}(\Phi) \cap \Sigma_1$ the frame for $\Sigma_1$ consisting of the drawing
\begin{equation*}
(bd(\Sigma_1), \psi^{-1}(V(\Phi) \cup V(K)) \cap \Sigma_1),
\end{equation*}
where for each edge $R$ of this drawing
